\chapter*{General Conclusion}
\addcontentsline{toc}{chapter}{General Conclusion}

This final-year project addressed the detection of series arc faults in residential electrical installations --- a safety-critical problem left open by conventional protection devices, and whose central difficulty is generalization: a detector must recognize an arc signature it has never seen, through a load it has never seen, without ever raising a false alarm. The answer proposed by this work is a \emph{detection system} rather than a lone classifier: several complementary experts observing a multi-cycle window, merged by a voting decision layer. This report presented the realization of its first and most fundamental pillar --- the single-cycle detector Arc-FaultNet --- together with the data and evaluation infrastructure on which every future component of the system will rest.

\medskip

The first contribution is a complete, reproducible \textbf{data engineering pipeline}: controlled arc experiments acquired on the Institut Jean Lamour bench (three channels at 1~MHz, varied loads and electrode materials over two campaigns), segmentation into individual 50~Hz cycles anchored on the mains voltage, automatic labeling using the arc voltage as a laboratory-only oracle with a calibrated three-zone rule, per-cycle normalization eliminating the amplitude shortcut, and anti-aliased decimation to 102.4~kHz justified by both the physics of the arc (signatures below 50~kHz) and the constraints of embedded deployment. The result is a versioned corpus of 10\,860 labeled cycles with per-sample metadata, plus a never-touched inter-campaign holdout.

\medskip

The second contribution is \textbf{Arc-FaultNet} itself: a $\approx$~315k-parameter dual-branch network in which a temporal branch reads four physics-informed channels (normalized current, first difference, Teager--Kaiser energy, sliding RMS), a spectral branch reads the STFT spectrogram through a learnable frequency gate with time-preserving pooling, and a bidirectional multi-head Q/K/V cross-attention aligns and fuses the two views before a regularized classification head. On the combined dataset the final model reaches \textbf{98.8\% accuracy and 98.7\% F1}, and the controlled comparison shows that the sequence-level cross-attention contributes $+0.74$ points of accuracy and $+1.85$ points of recall over the initial gated fusion --- recovering precisely the hard positives that matter for a safety device. This fusion revision itself illustrates the method followed throughout: the gated fusion was first chosen under an explicit small-data-regime hypothesis, the error analysis identified its structural limit (the temporal bottleneck), and the replacement was adopted with its trade-off stated in advance --- alignment capability and a 48k-parameter reduction gained, a quadratic attention cost and a data-hunger risk accepted and mitigated.

\medskip

The third contribution is methodological. The project advanced by \textbf{diagnosis rather than accumulation}: an initial 100\% score was identified as interpolation-regime overfitting; a capacity probe reframed the problem from model size to representation quality; a negative result (the inter-cycle residual channel) was traced to its structural cause and converted into a design principle; a flawed ablation study --- confounded variables, uncontrolled capacity, single-seed claims --- was detected, corrected, and replaced by a one-factor-at-a-time protocol that finally ranked the components: dual-branch structure first ($-9.6$ points when reduced to a single-branch CNN), spectral and temporal branches next ($-3.2$ and $-2.6$), cross-attention fourth ($-1.7$), with squeeze-and-excitation and the deep head established as stability tools (28--51\% variance reduction across seeds) rather than peak-performance components. Recording-level GroupKFold validation then quantified what a random split cannot: on entirely unseen experimental conditions, accuracy settles at $93.4\% \pm 8.1$, with high precision preserved ($97.2\%$) and the loss concentrated in per-cycle recall ($87.9\%$).

\medskip

These last numbers, together with the forensic analysis of the three residual false positives --- a start-up transient seen on two consecutive cycles, and a degraded contact plausibly micro-arcing --- establish the report's closing argument: the remaining errors are \emph{physically unsolvable at the single-cycle scale}. What distinguishes an inrush transient from an arc is not any property of one cycle, but persistence across cycles.

\medskip

The perspectives follow directly, and constitute the planned continuation of this work:

\begin{itemize}[itemsep=3pt]
    \item \textbf{The memory-based expert.} A sequence model with internal state --- a state space model (SSM) or a recurrent architecture --- consuming the stream of cycles (or of Arc-FaultNet's per-cycle embeddings) to recognize the temporal pattern of arc occurrence: persistence, intermittency, ignition/extinction dynamics over a window of up to eight cycles, as tolerated by IEC~62606;
    \item \textbf{The third expert and the voting layer.} A detector built on a different principle (e.g.\ descriptor-based statistics or inter-cycle residual analysis, where that information naturally belongs), chosen for decorrelated failure modes, and a voting rule ($k$-of-$n$ over the window, possibly weighted by per-expert confidence) tuned on deployment-level metrics: reaction time, missed-fault probability per event, nuisance trips per operating hour;
    \item \textbf{Data expansion toward the field.} New campaigns with masking loads (dimmers, switched-mode supplies), combined and parallel loads, and ultimately recordings from real domestic installations --- the regime where the attention mechanisms and the ensemble design are expected to show their full value;
    \item \textbf{Embedded deployment.} Quantization and export of the $\approx$~315k-parameter model to a microcontroller target, and measurement of the true per-cycle inference latency within the 20~ms budget.
\end{itemize}

\medskip

Beyond its technical outcome, this project was an exercise in engineering judgment: trusting no number before understanding its protocol, treating every failure as information, and designing not a model that scores well, but a system whose remaining errors are understood --- and already have an owner in the architecture of its next version.
